\scanpage{5}
% L1
Soient $\Gamma$ un groupe, $(\underline{\tau}_i)_{0\leq i\leq2}$ \deleted{\unsure{un}} syst. de générateurs de $\Gamma$, avec $\underline{\tau}_i^2=1$, $\underline{\tau}_0\underline{\tau}_2=\underline{\tau}_2\underline{\tau}_0$
% L2
On pose $\underline{\rho}_0=\underline{\tau}_0\underline{\tau}_1$, $\underline{\rho}_1=\deleted{\unsure{\tau_2\tau_1}}\;\underline{\tau}_1\underline{\tau}_2$.

% L3
\uline{Thm} \deleted{Soit} $S$ schéma, \deleted{\unsure{$G$ une forme de $GP(1)_S$}, \unsure{$G$ un groupe abstrait}}
% L4
\deleted{\unsure{eng.}} Considérons :

% L5
A) Catégorie $C$ des couples $(G,\varphi)$, où $G$ est une forme de
% L6
$\underline{GP}(1)_S$, $\varphi:\Gamma\longrightarrow G(S)$, satisfaisant les conditions
% L7
$(*)\ \forall s\in S$, $\rho_{0s}^{2}$ et $\rho_{1s}^{2}$ sont $\ne1_s$ [ NB on désigne par $\tau_0,\tau_1,\tau_2,\rho_0,\rho_1$ les
% L8
images de $\underline{\tau}_0,\underline{\tau}_1,\underline{\tau}_2,\underline{\rho}_0,\underline{\rho}_1$ dans $G(S)$ ] et $\rho_{0s}$ et $\rho_{1s}$ ne commutent pas

% L9-13
\deleted{$(\Leftrightarrow)$ b) $\forall s\in S$, l'image $\Gamma_s$ de $\Gamma$ dans $G(s)$ [$\simeq\Gamma/N_s$, où $N_s$ est le noyau de $\Gamma\to G(s)$] \uline{n'est} pas isomorphe à un groupe diédral $\mathbb D_n$ \unsure{($n\geq2$, fini)|($n\geq3$)} $\geq3$ -- le cas $n=1$ ou $2$ est déjà exclu par a), qui signifie que $\tau_{0s}$ et $\tau_{1s}$ ne commutent pas, ni $\tau_{1s}$ et $\tau_{2s}$ -- donc $\Gamma_s$ n'est pas commutatif].}
% L10a
\inserted{\deleted{(ou : un produit $\mathbb D_\infty\times\mathbb Z_{/2}$ \unsure{$\text{ou }\mathbb Z_{/2}\times\mathbb Z_{/2}$})}}
% L11a
\margin{\deleted{On peut remplacer b) par b') $\rho_0$ et $\rho_1$ ne commutent pas (ou encore $\rho_0$ et $\sigma=\rho_0\rho_1=\tau_0\tau_2$ ne comm. pas (NB $\sigma^2=1$)}}

% L14
Les flèches de $C$ sont les isomorphismes.

% L15
B) La catégorie discrète $C'$ définie par l'ens des couples
% L16
$(\deleted{\unsure{\beta_0(\varphi)|\alpha_0(\varphi)}}\alpha_0\unsure{(\varphi)},\deleted{\unsure{\beta_1(\varphi)|\alpha_1(\varphi)}}\alpha_1\unsure{(\varphi)})\in\Gamma(S,\underline O_S)^2$, satisfaisant les conditions
% L17
\[\begin{aligned}&\deleted{\unsure{\beta_0,\beta_1,\beta_0-4,\beta_0+\beta_1}\ \text{sont}}\\
% L18
&\alpha_0+2,\ \alpha_1+\unsure{2|1},\quad\alpha_0+\alpha_1\quad\text{sont inversibles.}\end{aligned}\tag*{(*)}\]

% L19
Pour tout objet \deleted{\unsure{de}} $(G,\varphi)$ de $C$, soit
% L20
\[\alpha_0=A(\deleted{\unsure{\varphi}}\rho_0),\quad\alpha_1=A(\varphi_1)\]
% L21
où
% L21a
\[A:G\mathbin{-}\mathbb E^1_S\]
% L22
est l'invariant. \deleted{On trouve} Ceci posé

% L23
1) $(\alpha_0,\alpha_1)$ ainsi associé à $(G,\varphi)$ satisfait $(*)$

% L24
2) Le foncteur
% L25
\[ C\longrightarrow C' \]
% L26
défini ainsi par $(G,\varphi)\longmapsto(\alpha_0(\varphi),\alpha_1(\varphi))$ est \uline{\unsure{pl\textsuperscript{t}.|pl.} fidèle},
% L27
et c'est une \unsure{$\left.\right\}$} équivalence de catégories \inserted{si $\Gamma=\{\tau_0,\tau_1,\tau_2\mid\tau_0^2=\tau_1^2=\tau_2^2=(\tau_0\tau_2)^2=1\}=G_2$.} En d'autres termes

% L28
a) Le groupoïde $C$ est \uline{rigide}, i.e. $\forall (G,\varphi)$, le centralisateur de
% L29
$\varphi$ dans $G(S)$ est $\{1\}$.

% L30
b) Deux objets $(G,\varphi)$ et $(G',\varphi')$ de $C$ sont isom ssi
% L31
\[\alpha_0(\varphi)=\alpha_0(\varphi'),\quad\alpha_1(\varphi)=\alpha_1(\varphi')\]

% L32
c) \inserted{Si $\deleted{\unsure{G}}\Gamma=G_2$} $\forall$ couple $(\alpha_0,\alpha_1)$ \deleted{\unsure{$\in A$}} de sections de $\underline O_S$ satisfaisant $(*),$
% L33
$\exists\ (G,\varphi)$ \deleted{tel q} $\mathrm{Ob}\ C$ tel que $\alpha_0(\varphi)=\alpha_0$, $\alpha_1=\alpha_1(\varphi)$ [et
% L34
ledit objet de $C$ est défini à isom. unique près grâce à a)]
% archive sheet number: 4
